
\begin{table}[t]
\centering
\caption{Crash resolution rate (CRR) for different tools on the \kbenchsyz{} benchmark ($200$ bugs). LLMs used by the tools are in parentheses. $^*$Results are from \citet{crashfixer}, out of $279$ bugs.
}
\label{tab:rq1-main}
\vspace*{-2mm}
\scalebox{0.95}{
\begin{tabular}{ccccc}
\cmidrule{1-5}
\textbf{Setting} & \textbf{Tool} & \textbf{Max calls}& \textbf{P@k} & \textbf{CRR (\%)} \\
\cmidrule{1-5}
\multirow{3}{*}{\rotatebox[origin=c]{0}{\centering Assisted}} 
  & GPT-4o           & $1$ & P@$5$ & $36.00$ \\
  & o1                       & $1$ & P@$5$  & $51.00$ \\
  & CrashFixer~(Gemini 1.5 Pro-002)$^*$ & $\geq 4$ & P@$16$ & $49.22^*$ \\
  & CrashFixer~(Gemini 2.5 Pro) & $\geq 4$ & P@$16$ & $\mathbf{70.00}$ \\
\cmidrule{1-5}
\multirow{2}{*}{\rotatebox[origin=c]{0}{\centering Stack context}} 
  & GPT-4o & $1$ & P@$5$  & $29.50$ \\
  & o1 & $1$ & P@$5$  & $\mathbf{40.00}$ \\
\cmidrule{1-5}
\multirow{3}{*}{\rotatebox[origin=c]{0}{\centering Unassisted}} 
  & Agentless (GPT-4o) & $4$ & P@$5$  & $31.00$ \\
  & SWE-agent (GPT-4o) & $15$ & P@$5$  & $31.50$ \\
  & \textbf{\cresearcher{} (GPT-4o)} & $15$ & P@$5$  & $48.00$ \\
  & \textbf{\cresearcher{} (GPT-4o + o1)} & $15$ & P@$5$  & $58.00$ \\
  & \textbf{\cresearcher{} (Gemini 2.5-Flash)} & $15$ & P@$5$  & $\mathbf{67.00}$ \\
\cmidrule{1-5}
\multirow{4}{*}{\rotatebox[origin=c]{0}{\centering \makecell{Unassisted \\ + Scaled} }} 
  & SWE-agent (GPT-4o) & $30$ & P@$5$  & $32.00$ \\
  & \textbf{\cresearcher{} (GPT-4o)} & $30$ & P@$5$  & $47.50$ \\
  & SWE-agent (GPT-4o) & $15$ & P@$10$ & $37.50$ \\
  & \textbf{\cresearcher{} (GPT-4o)} & $15$ & P@$10$  & $\mathbf{54.00}$ \\
\cmidrule{1-5}
\end{tabular}
}
\end{table}

\subsection{RQ1: How effective are different tools at resolving Linux kernel crashes?}%
\label{sec:rq1}

Our main results are presented in Table~\ref{tab:rq1-main}, organized by setting, namely, assisted, stack context, unassisted, and unassisted + test-time scaled (Section~\ref{sec:experiments}).

\noindent\textbf{1) The assisted setting baselines show strong performance, but under unrealistic assumptions.}
Given the ground-truth buggy files, LLMs like GPT-4o (CRR of $36\%$) are quite capable of resolving crashes.
A reasoning model like o1 is significantly better (CRR of $51\%$).
Adding iterative feedback from the compiler and the crash reproduction setup, CrashFixer achieves $49.22\%$ CRR with an older Gemini model and $70\%$ CRR with a newer Gemini reasoning model
using P@$16$ with at least $4$ \maxcalls{}.

\noindent\textbf{2) The stack context setting reveals the difficulty of the practical unassisted setting.}
The assumption that an oracle can tell us exactly which files need to be edited is impractical.
The gap between the assisted (idealistic) setting and the practical unassisted setting is highlighted by the simple but effective stack context setting
where models are given the contents of all the files mentioned in the crash report (truncated to the context length limit).
This is a strong baseline because all the ground-truth buggy files are present in the crash report for {$74.50\%$} crashes in our dataset.
o1 achieves a CRR of $40\%$, which is impressive, but $11\%$ lower than its performance in the assisted setting.

\noindent\textbf{3) \cresearcher{} consistently outperforms baselines in the unassisted setting.}
Fixing the LLM as GPT-4o, \cresearcher{} achieves a CRR of $48\%$, significantly outperforming the SWE-agent and Agentless baselines, both the stack context baselines, and even the assisted GPT-4o baseline.
This indicates that the context gathered by \cresearcher{} is much more effective than giving file contents based on the crash report, or using multi-step hierarchical localization (as done by Agentless), and is even better than directly giving all the contents of the files to be edited.
The context gathered by GPT-4o during \analysisphase{} can be better utilized by the reasoning model o1 during \synthesisphase{}, increasing \cresearcher{}'s CRR to $58\%$.
This further increases to $67\%$ using the newer (albeit not as advanced as Gemini 2.5-Pro) model, Gemini 2.5-Flash, for both phases, indicating that \cresearcher{}'s design is not tied to a specific LLM family.

\noindent\textbf{4) Increasing the number of trajectories helps, while making them longer does not.}
We examine how scaling the total inference budget (\maxcalls{} $\times$ num trajectories $k$) impacts the performance of \cresearcher{} and SWE-agent.
Doubling the \maxcalls{} budget, i.e., making the trajectories of the agents longer, has a negligible effect on the CRR, whereas
increasing the number of trajectories sampled improves SWE-agent's CRR to $37.50\%$ and \cresearcher{} (GPT-4o)'s CRR to $54.00\%$.
